\chapter{Literature review}
\section{Chapter overview}
Describe how this chapter reviews concepts, technologies, existing work, and evaluation practices.

\section{Concept map}
Insert a conceptual map figure and provide a brief explanation of concept relationships.

\begin{figure}[H]
    \centering
    \fbox{\rule{0pt}{2in}\rule{0.9\linewidth}{0pt}}
    \caption{Concept map of the research domain (placeholder)}
    \label{fig:concept-map}
\end{figure}

\section{Problem domain review}
\subsection{Introduction to the domain}
Define the overarching domain and its key terms.

\subsection{Key challenges and sub-areas}
Review major thematic challenges relevant to your problem.

\subsection{Existing frameworks, standards, and domain approaches}
Summarize applicable standards, regulations, and best-practice frameworks.

\subsection{Proposed architecture positioning}
Position your high-level architecture against the reviewed domain context.

\begin{figure}[H]
    \centering
    \fbox{\rule{0pt}{2in}\rule{0.9\linewidth}{0pt}}
    \caption{Proposed conceptual architecture (placeholder)}
    \label{fig:proposed-architecture}
\end{figure}

\section{Technological review}
\subsection{Data acquisition and preprocessing}
Critically review data handling techniques used in comparable work.

\subsection{Model or system architectures}
Compare relevant architecture choices and trade-offs.

\subsection{Tuning and optimization approaches}
Review methods used to optimize performance and robustness.

\subsection{Evaluation metrics used in prior work}
Evaluate suitability of common metrics for your project context.

\section{Existing work}
\subsection{Thematic grouping}
Group studies by themes and compare their findings.

\subsection{Cross-theme linking}
Explain overlaps and differences across themes and methods.

\section{Evaluation and benchmarking}
Critically analyze datasets, protocols, metrics, and evaluation limitations.

\section{Chapter summary}
Synthesize key gaps and connect them to Chapter 3.
